% !TeX root = presentation.tex
% !TeX spellcheck = de_DE

\documentclass[presentation.tex]{subfiles}

\begin{document}
	\frame{\titlepage}
    
    \begin{frame}
        \frametitle{Überblick}
        \begin{itemize}
            \item Einführung Lorenz Attraktor
            \item Analyse der Gleichungen
            \item Durchspielen Szenarios
            \item Outro
        \end{itemize}
    \end{frame}
    
    \begin{frame}
    	\frametitle{Einführung Lorenz Attraktor}
    	\framesubtitle{Butterfly Effekt}
		\textit{"Predictability: Does the Flap of a Butterfly's Wings in Brazil Set Off a Tornado in Texas?"}, \\
		Paper Edward N. Lorenz, 1979
    \end{frame}
    
    \begin{frame}
    	\frametitle{Einführung Lorenz Attraktor}
    	\framesubtitle{Eigenschaften}
    	\begin{itemize}
    		\item Modell für Zirkulation der Atmosphäre auf der Erde
    		\item Zeigt, dass das Wetter chaotisches Verhalten hat
    		\item Näherung, keine verlässliche Vorhersage
    	\end{itemize}
    \end{frame}
    
    \begin{frame}
		\frametitle{Einführung Lorenz Attraktor}
		\framesubtitle{Was ist ein Attraktor?}
		
    \end{frame}
    
	\begin{frame}{Experiment zum Lorenz Attraktor}
		\begin{columns}[c]
			\begin{column}{.5\textwidth}
				\url{https://lorenz.olidias.ch}
			\end{column}
			\begin{column}{.5\textwidth}
				\includegraphics[width=\textwidth]{olidiasQRCode}
			\end{column}
		\end{columns}
	\end{frame}
\end{document}